\section{Complete native source-entry conjugation in every rank}
\label{sec:r10-native-source-diagonal-shear-conjugation}

\paragraph{Status and scope.}
\textbf{PROVED THIS ROUND:} the prescribed full SHT elementary-shear
pair and the prescribed SHT source-entry unit pair have an explicit
strict second witness in every rank.  The shear coefficient is
arbitrary analytic, with zeros and unbounded values allowed.  All
valuation walls, mixed coefficients and the original metric remain
in the actual formulas.  Combining this result with MX gives
conjugation by any supplied finite diagonal-unit path, without a
common bounded diagonal-frame hypothesis.

This resolves the native source-entry gap left in MX and RSC.
It uses a different direct six-edge witness; it does not identify
the all-entry native scalar path with the CNA scalar-reference
pair.  That latter comparison is still open.  The older R5 SD
rank-two source relation for its specified ET/UD rule is a frozen
different-rule input, not recounted here.  The root and diagonal
agents are construction contributors to this section.

\subsection{The original-field six-edge loop}

Let $N$ be an ordinary analytic open over the original $K$, let
$f$ be an ordered $r$-tuple of analytic sections of a line bundle
with its fixed metric, and put
\[
        r\geq2,\quad s=r-1,\quad W=N\setminus|f=0|,
                 R=c_1(\overline L),\quad D=d'd''.
 \tag{NSC1}
\]
Fix $i\ne j$, an analytic coefficient $t\in\mathcal O(N)$ and
an analytic unit $\lambda\in\mathcal O(N)^\times$.
Write $E(t)=I+t e_{ij}$ and
\[
       g=E(t)f,\qquad k=D_j(\lambda)f,\qquad
       h=D_j(\lambda)g=E(t/\lambda)k .
 \tag{NSC2}
\]
Thus the exact matrix relation is
$D_j(\lambda)E_{ij}(t)=E_{ij}(t/\lambda)D_j(\lambda)$.
All these tuples have the same zero locus as $f$.
Let $\mathcal P_{ij,t}(f)$ be the native SHT shear pair,
including SHT's actual zero gluing, and let
$p_{j,\lambda}(f)$ be its prescribed one-entry unit pair.
We shall construct an actual triple $Q_{\rm src}\in L_2^s(W)$
with
\[
 \mathcal P_{ij,t}(f)+p_{j,\lambda}(g)
       -p_{j,\lambda}(f)-\mathcal P_{ij,t/\lambda}(k)
                            =\partial_2Q_{\rm src}.
 \tag{NSC3}
\]
The actual pair/triple spaces, full differentials and negative-degree
conventions are CPT1; in particular
$\partial_2(A,B,M)=(d'A+d''M,d''B+d'M)$.

First work on the ordinary open $N_t=\{t\ne0\}$, and put
$a=v(\lambda)$, $d=v(t)$.  Denote the following genuine tuple
potentials by capitals:
\[
 \begin{gathered}
 P=\Phi_f,\quad Q=\Phi_g,\quad H=\Phi_h,\quad K=\Phi_k,\\
 U=\Phi_{D_i(t)^{-1}f}=\Phi_{D_i(t)^{-1}g},\\
 V=\Phi_{D_i(t/\lambda)^{-1}h}
                     =\Phi_{D_i(t/\lambda)^{-1}k}.
 \end{gathered}
 \tag{NSC4}
\]
Each equality in the second and third lines is the literal
norm equality supplied by $E_{ij}(1)$ and its inverse.
The full defect on the left of (NSC3) is the sum of the
six native SHT unit edges
\[
 \begin{array}{c|c}
 \text{potential edge}&\text{unit weight}\\ \hline
 P\longrightarrow U&-d\\
 U\longrightarrow Q&d\\
 Q\longrightarrow H&a\\
 H\longrightarrow V&a-d\\
 V\longrightarrow K&d-a\\
 K\longrightarrow P&-a .
 \end{array}
 \tag{NSC5}
\]
The middle norm-isometric tuple steps in SHT have zero pair.
The weights sum to zero and every endpoint expression telescopes.
We retain the actual potential $V$ in every rank.  In general
$V$ is not $U+a$ when passive entries are present.

\subsection{A complete degree-one relation on the native triangle}

For either one of the two full first differentials separately
write
\[
  \mathfrak s(u,v)=u\,dv-v\,du,\qquad
  S(p,q,c)=\mathfrak s(p,q)+\mathfrak s(c_+,q)
                                      +\mathfrak s(c_-,p),
  \quad c_+=\max(c,0),\quad c_-=\min(c,0).
 \tag{NSC6}
\]
The letter $d$ in $\mathfrak s$ denotes that differential only
in this subsection; the scalar $d=v(t)$ remains an argument
when it occurs in $\mathfrak s(d,a)$.  Reversing an edge gives
$S(q,p,-c)=-S(p,q,c)$.  A common actual scalar shift $b$ gives
\[
 S(p+b,q+b,c)-S(p,q,c)
               =\mathfrak s(c-(q-p),b).
 \tag{NSC7}
\]
Therefore the sum around (NSC5) is invariant under every
common scalar shift, with no assumption $Db=0$.

\begin{lemma}[The entire six-edge degree-one residual vanishes]
If all six tuple norms are determined by entries $i,j$, then
\[
 \begin{split}
 &S(P,U,-d)+S(U,Q,d)+S(Q,H,a)\\
 &\qquad+S(H,V,a-d)+S(V,K,d-a)+S(K,P,-a)=0 .
 \end{split}
 \tag{NSC8}
\]
This is an equality of full actual first forms in each colour,
including every unit-valuation wall and all zero-entry cases.
\end{lemma}

\begin{proof}
First suppose $f_j,f_i,f_i+t f_j$ are nonzero on the chart.
Subtract the common potential of $f_j$, which is permitted
by (NSC7), and set
\[
 x=v(f_i/f_j),\quad y=v((f_i+t f_j)/f_j),\quad
 M=\min(x,d)=\min(y,d),\quad u=M-d .
 \tag{NSC9}
\]
The equality of the two minima is the actual norm equality
$|f_i/t| \mathbin{\rm max}|f_j|
 =|(f_i+t f_j)/t| \mathbin{\rm max}|f_j|$.
Moreover the minimum of $x,d,y$ is attained at least
twice, by the actual analytic relation
$(f_i+t f_j)/f_j=f_i/f_j+t$.
The six normalized potentials are
\[
 p=\min(x,0),\quad q=\min(y,0),\quad
 k=\min(x,a),\quad h=\min(y,a),\quad u,\quad v=u+a .
 \tag{NSC10}
\]
For $X=x$ or $y$ write $p_X=\min(X,0)$, $k_X=\min(X,a)$
and put
\[
 F(X)=S(p_X,u,-d)+S(u+a,k_X,d-a)+S(k_X,p_X,-a).
 \tag{NSC11}
\]
By exact edge reversal, the left side of (NSC8) is
$F(x)-F(y)$.

On the part $M<d$ of the native triangle, $x=y=M$,
so the difference is zero with its entire first jet.
On the part $M=d$, one has $u=0$ and $x,y\geq d$.
For any $X\geq d$, direct expansion of (NSC6) gives
\[
 \begin{split}
 F(X)={}&\mathfrak s(\max(a,d)-a_+,k_X)\\
        &+\mathfrak s(k_X-a_--d_+,p_X)
                           +\mathfrak s(\min(a,d),a).
 \end{split}
 \tag{NSC12}
\]
We prove $F(X)=\mathfrak s(d,a)$ on this entire part,
not just away from intercept walls.  Put $L=\max(a,d)$.
Since $\min(a,d)+L=a+d$, subtracting the desired value
from (NSC12) gives
\[
 \mathfrak s(L-a_+,k_X)
       +\mathfrak s(k_X-a_--d_+,p_X)-\mathfrak s(L,a).
 \tag{NSC13}
\]
There are four exhaustive closed position ranges.
\begin{enumerate}
\item If $X\geq\max(a,0)$, then $p_X=0$, $k_X=a$.
      Expression (NSC13) is $-\mathfrak s(a_+,a)=0$.
\item If $X\leq\min(a,0)$, the condition $d\leq X$ gives
      $L=a$ and $d_+=0$, while $p_X=k_X=X$.
      Thus (NSC13) is
      $\mathfrak s(a_-,X)+\mathfrak s(X-a_-,X)=0$.
\item If $0\leq X\leq a$, then $a_+=a$, $a_-=0$,
      $L=a$, $p_X=0$, $k_X=X$.  Every term in (NSC13)
      is zero.
\item If $a\leq X\leq0$, then $a_+=0$, $a_-=a$,
      $d_+=0$, $p_X=X$, $k_X=a$.  The first and last
      terms cancel and the middle one is zero.
\end{enumerate}
All equalities used here are equalities of the restricted
scalar functions on the relevant polyhedral cells and hence
of their full tangential first jets.  In particular
$\mathfrak s(a_+,a)=0$ also on the face $a=0$.
The four formulas agree on their common faces.
Refine the whole actual unit presentation by all these
minima and absolute-value walls.  The first-form coefficients
then vanish on each carrier with its tangent restrictions.
There is no positive-codimension first-colour or second-colour
Dirac term hidden in a first derivative of a continuous
piecewise affine scalar.  Thus this proves an actual
degree-one relation on every further refinement, not a
test on separated probe products.

Zero entries require ordinary neighborhoods, not a use of
an undefined valuation.  At a point with $f_j\ne0$ and
$f_i=0$, choose an ordinary neighborhood on which
\[
 |f_i|<\min(|f_j|,|\lambda f_j|,|t f_j|).
 \tag{NSC14}
\]
The normalized $p,k,u$ are then literally $0,a,0$, and
$y=d$; these are exactly the stable $x=+\infty$ values
in the preceding calculation.  They prove the same
identity without evaluating $v(0)$.
At a point with $f_j\ne0$ and $f_i+t f_j=0$, apply
the analogous inequality to $f_i+t f_j$; then
$q,h,u$ are $0,a,0$ and $x=d$.
All these inequalities compare sections of the same line
bundle, so are independent of its local frame.

Finally, if $f_j=0$, then $f_i\ne0$ under the two-entry
hypothesis.  On an ordinary neighborhood one has
$|f_i|>\max(|f_j|,|t f_j|,|\lambda f_j|)$.
With $b=\Phi_{f_i}$ the six potentials are
$P=Q=H=K=b$, $U=b-d$, $V=b+a-d$.
Edges one and two, three and six, and four and five
are three exact inverse pairs.  Their residuals cancel
literally.  The units $t,\lambda$ are finite and nonzero
at each point of $N_t$, so these neighborhoods exist.
This covers every zero-entry case.  Adding back the
arbitrary common scalar or metric by (NSC7) proves (NSC8).
\end{proof}

The lemma is an identity of actual first forms; it is
stronger than a statement that the resulting pair is
closed.  In the CS notation it says
$\sum_e S_{e,1}=0$ with both colour signs retained.

\subsection{The finite edge triples and their reversal}

We use the proved full edge identity CS10 with all its
wall terms.  To specify the actual triple completely,
for an edge $(p,q,c)$ and reference $R_0$ put
\[
 \begin{gathered}
 C_p=R_0-Dp,\quad C_q=R_0-Dq,\quad z=|c|,\\
 H_k=h_{k-1}(R_0,C_p,C_q),\quad
 J_k=h_{k-1}(C_p,C_q),\\
 K_{p,k}=h_{k-1}(R_0,C_p),\quad
 K_{q,k}=h_{k-1}(R_0,C_q),\\
 M_k=-\tfrac14\bigl[(q^2-p^2)H_k+z(q-p)J_k
                               +c(pK_{p,k}+qK_{q,k})\bigr].
 \end{gathered}
 \tag{NSC15}
\]
Here $h_j$ is the complete homogeneous polynomial in
all displayed variables, with $h_j=0$ for $j<0$.
Let $Q_k^{\rm tet}(R_0;z,p,q)$ be the explicit triple
CS8--CS9 for the four potentials $(0,z,p,q)$.
In particular it uses
$h_{k-2}(R_0,R_0-Dz,C_p,C_q)$ and the full three
alternating first-jet numerators in CS9.
Set $Q_1^{\rm tet}=Q_0^{\rm tet}=0$ and define
\[
 \mathcal K_k(p,q,c;R_0)
     =(0,0,M_k)
          +\tfrac12\bigl(Q_k^{\rm tet}-R_0Q_{k-1}^{\rm tet}\bigr).
 \tag{NSC16}
\]
This is a specified finite actual triple, not a chosen
primitive modulo a quotient.  The frozen-in-this-round
identity and its two residual support rules are
\[
 \begin{gathered}
 p_e^{[k]}=\partial_2\mathcal K_k-\tfrac12 S_{e,k},\\
 S_{e,k}=R_0S_{e,k-1}\quad(k\geq2)
       \quad\text{on an open where only the two changed entries
       determine all relevant norms}.
 \end{gathered}
 \tag{NSC17}
\]
The second line is the edgewise assertion CS17, whose
proof CS19 uses only the two-entry condition for that
edge, not the existence of a four-edge square.
It retains the coefficient identities on each carrier
before multiplying by any other curvature factors.
It never asserts that mixed products $C_pC_q$ vanish.

The edge triple has literal antisymmetry:
\[
          \mathcal K_k(q,p,-c;R_0)
                         =-\mathcal K_k(p,q,c;R_0).
 \tag{NSC18}
\]
Indeed every term of $M_k$ changes sign; $z$ and
all symmetric curvature polynomials stay fixed.
The three numerators of $Q_k^{\rm tet}$ change sign
when $p,q$ are interchanged, while its full
$R_0-D|c|$ factor is unchanged.  This verifies
(NSC18) through every absolute-value wall.

\subsection{One passive cap and the complete all-rank witness}

In rank two, use (NSC8) and the first line of (NSC17)
at $s=1$.  On $W\cap N_t$ the explicit scalar triple
\[
        Q^{(2)}=\sum_{e=1}^{6}(0,0,M_{e,1})
 \tag{NSC19}
\]
has boundary equal to the full six-edge defect.
The potentials in this formula are the original
$P,U,Q,H,V,K$; the proof of (NSC8) allowed their
arbitrary common metric term.

Now let $r\geq3$.  Entries other than $i,j$ form the
same passive tuple at every step.  Write its norm
potential as $\psi$, temporarily $+\infty$ where
that tuple vanishes.  On $W\cap N_t$ put
\[
 \begin{gathered}
 \kappa_0=\sum_{e=1}^6|c_e|
              =2|d|+2|a|+2|d-a|,\qquad \varepsilon>0,\\
 \Psi=\min(\psi,P+\kappa_0+\varepsilon),\quad
 R_*=R-D\Psi,\quad p_v=\Phi_v-\Psi .
 \end{gathered}
 \tag{NSC20}
\]
This is one actual scalar on the entire stated domain.
On the passive-nonzero locus it has the finite
piecewise smooth presentation of CS24.  At a passive-zero
point the ratio of the passive norm to the full norm
vanishes, while $e^{-\kappa_0-\varepsilon}$ is positive
and continuous.  On an ordinary neighborhood the cap
therefore equals $P+\kappa_0+\varepsilon$ literally.
These equalities give actual sheaf gluing across the
passive-zero locus.  No analytic element with valuation
$\varepsilon$ is needed; all real intercepts are allowed.
The original common metric occurs once and is retained.

A one-entry unit change of weight $c_e$ changes a norm
potential by at most $|c_e|$.  Following (NSC5), including
its norm-isometric intermediate tuples, gives
\[
          \Phi_v\leq P+\kappa_0
                    \quad(v=P,U,Q,H,V,K).
 \tag{NSC21}
\]
The two ordinary analytic opens
\[
 \begin{split}
 W_{\rm pas}&=\{\psi<P+\kappa_0+\varepsilon\},\\
 W_{\rm two}&=\{\psi>P+\kappa_0\}
 \end{split}
 \tag{NSC22}
\]
cover $W\cap N_t$.  Their strict norm-ratio descriptions
also show that $W_{\rm two}$ contains every passive-zero
point.

On $W_{\rm pas}$, the cap is the passive potential.
Its projective curvature $R_*$ is pulled back by the
actual original-field map to $\mathbf P^{r-3}$ defined
by the $r-2$ passive entries.  Hence $R_*^{r-2}=0$.
The proof of CS21--CS22 is edgewise and applies to
each of the six unit edges, giving $S_{e,s}=0$.
It distinguishes a changed active entry, an active
passive entry, and the unique other active entry;
all its full carrier and first-jet arguments remain
valid here.  For $r=3$ the passive target is
$\mathbf P^0$ and $R_*=0$, so this endpoint is included.

On $W_{\rm two}$, (NSC21) makes the passive entries
strictly inactive at every vertex and intermediate
tuple.  Only here may the rank-two normalized
identification $v=u+a$ from (NSC10) be used.
The degree-one identity (NSC8), invariant under the
common shift by $-\Psi$, and the edgewise recurrence
(NSC17) yield
\[
              \sum_eS_{e,s}
                    =R_*^{s-1}\sum_eS_{e,1}=0 .
 \tag{NSC23}
\]
The recurrence allows any common closed reference
$R_*$.  Thus it retains the full cap curvature,
including all derivatives of $\kappa_0$ and all cap
transition walls.  The transition itself lies in
$W_{\rm two}$ because $\varepsilon>0$.
Both sides of (NSC23) are global actual formulas;
the two-open argument verifies their equality and
does not extend a locally chosen triple.

For clarity specify the complete metric triple too.
Use the original six potentials and weights in MT5--MT6
to form $f_k$ and $\Gamma_j$, put $b=-\Psi$, and let
\[
 \begin{gathered}
 T_\ell(R_*,R)=\sum_{u=0}^{\ell}R_*^uR^{\ell-u},\\
 Q_{\rm met}=
  \sum_{k=1}^{s-1}T_{s-k-1}(R_*,R)H_b(f_k)
       -(0,0,b\sum_{k=1}^{s}\Gamma_{k-1}R_*^{s-k}),\\
 H_b(f_k)=
       (d''b\wedge f_k',-d'b\wedge f_k'',b\,d'f_k').
 \end{gathered}
 \tag{NSC24}
\]
All six weights sum to zero and all endpoints telescope,
so MT16 applies to this entire loop:
the normalized loop minus the original loop is
$\partial_2Q_{\rm met}$.  In particular MT's full
nonclosed coefficient remainder has not been omitted.

Define on $W\cap N_t$, for $r\geq3$,
\[
 Q^{(r)}=
    \sum_{e=1}^{6}\mathcal K_s
           (p_{\mathrm{start}(e)},p_{\mathrm{end}(e)},c_e;R_*)
                                      -Q_{\rm met}.
 \tag{NSC25}
\]
The sum of the first lines of (NSC17), followed by
(NSC23) and the metric identity, proves
$\partial_2Q^{(r)}$ equals the full original defect.
There is no remaining extending pair.  Every term
in (NSC25) is actual on the entire $W\cap N_t$.

\subsection{Literal zero gluing across the shear coefficient}

The ordinary open
\[
                 O=\{|t|<1,\ |t/\lambda|<1\}\subset N
 \tag{NSC26}
\]
contains every zero of $t$.  On it the shears $E(t)$
and $E(t/\lambda)$ and their inverses are norm
isometries.  On $O\cap N_t$ this gives
$Q=P$ and $H=K$.  Consequently the six potential
edges in (NSC5) are three exact inverse pairs:
one with two, three with six, and four with five.

For rank two, (NSC18) makes $Q^{(2)}=0$.
For every higher rank, the same single cap $\Psi$
is used for the entire loop, so the normalized edges
are still exact inverse pairs.  Thus their
$\mathcal K_s$ triples sum to zero by (NSC18).
The entire sums $f_k$ and $\Gamma_j$ in MT5--MT6
also vanish by exact edge reversal.  Formula (NSC24)
then gives $Q_{\rm met}=0$, regardless of the cap's
value or its derivatives.  We have proved
\[
                    Q^{(r)}=0
                         \quad\text{on }W\cap O\cap N_t
                                  \quad\text{for every }r\geq2.
 \tag{NSC27}
\]
This is literal equality of actual triples and all
mixed coefficients, not equality of their boundaries.

Glue $Q^{(r)}$ on $W\cap N_t$ with the zero triple
on $W\cap O$.  These ordinary opens cover $W$.
On $O$ both native shear pairs in (NSC3) are zero,
and the two diagonal pairs have equal input potentials
$P=Q$, $K=H$, so their difference is zero too.
The glued actual triple therefore satisfies (NSC3)
everywhere, including components on which $t$ vanishes
identically.  The walls $|t|=1$, $|t/\lambda|=1$ and
every weight-zero wall lie in $N_t$ and remain in
the full formulas there.  No singular valuation
is extended through $t=0$.

\begin{theorem}[Complete native source-entry conjugation]
For every $r\geq2$, every analytic $t$ and every
analytic unit $\lambda$ on the stated ordinary open,
(NSC19), (NSC25) and the literal zero gluing (NSC27)
construct an actual $Q_{\rm src}$ satisfying (NSC3).
The field, absolute value, common metric and tuple
zero locus are unchanged.  No boundedness, valuation
order, analytic root or extra growth condition is required.
\end{theorem}

\subsection{Finite diagonal conjugation and the original-support bridge}

Let $\mathcal D$ be a supplied finite path of one-entry
analytic-unit diagonal factors with endpoint matrix
$D_0=\operatorname{diag}(\lambda_1,\ldots,\lambda_r)$.
Its prescribed pair $p_{\mathcal D}(f)$ is the literal
sum of its native SHT edges in the supplied order.
The exact matrix identity is
$D_0E_{ij}(t)=E_{ij}((\lambda_i/\lambda_j)t)D_0$.
Cross a shear past the diagonal factors successively.
For a factor in entry $i$, use MX2; for one in entry
$j$, use (NSC3); for every other entry, use MX8.
At each step use the actual intermediate tuple and
the then-current analytic shear coefficient.
Every coefficient remains analytic and may have zeros.
Summing the finitely many signed actual triples gives
\[
 \mathcal P_{ij,t}(f)+p_{\mathcal D}(E_{ij}(t)f)
       -p_{\mathcal D}(f)
       -\mathcal P_{ij,(\lambda_i/\lambda_j)t}(D_0f)
                              =\partial_2Q_{\mathcal D,E;f}.
 \tag{NSC28}
\]
Common intermediate primitive expressions cancel literally.
No common bounded diagonal frame or global dominant
unit is needed.  This is for the supplied diagonal
path; it does not assert an unspecified canonical
native primitive for every diagonal matrix.
The same finite telescoping compares a fixed finite
shear word with its coefficientwise conjugate.

For the exact-support consequence let $X/K$ be smooth
projective of dimension $n$ and keep the original
regular CI difference $Z=[Y_1]-[Y_0]$, with
the same line bundle and metric, and put
$V=X^{\rm an}\setminus(|Y_0|\cup|Y_1|)$,
$U=X^{\rm an}\setminus|Z|$.  Put
$C_U=(|Y_0|\cup|Y_1|)\cap U$ and supply an ordinary
neighborhood $C_U\subset N\subset U$ with an ordinary
cover $N=N_0\cup N_1$.  Suppose it has local actual descriptions
$G_1-G_0=\partial_1p_\nu+E_\nu$,
and their whole-overlap discrepancy is a specified
finite sum of the native conjugation relations above.
Then the corresponding finite sum $q=(A,B,M)$ is
an actual strict witness.  For a subordinate ordinary
partition $\theta_0+\theta_1=1$, form
\[
 \begin{split}
 \alpha={}&\theta_0\alpha_0+\theta_1\alpha_1
                      +d'\theta_1\wedge A+d''\theta_1\wedge M,\\
 \beta={}&\theta_0\beta_0+\theta_1\beta_1
                      +d''\theta_1\wedge B+d'\theta_1\wedge M,\\
 E_N={}&\theta_0E_0+\theta_1E_1 .
 \end{split}
 \tag{NSC29}
\]
The complete RAG calculation gives
$G_1-G_0=d'\alpha+d''\beta+E_N$; both mixed
$D\theta_1\wedge M$ terms are retained.
Choose the RSP cutoff $\chi$ on exact $U$, equal
to one near $C_U$ and with closed $U$-support contained
in $N$.  On $V$ set
\[
 \begin{gathered}
 g_V=G_1-G_0-d'\operatorname{ext}(\chi\alpha)
                    -d''\operatorname{ext}(\chi\beta),\\
 g_V=E_N\ \text{near the cancelled part},\qquad
 g_Z\in B^{r-1,r-1}(U),\qquad
          \int_Ug_Z\wedge D\eta=-\int_Z\eta
                   \quad(\eta\in B_c^{n-r,n-r}(X^{\rm an})).
 \end{gathered}
 \tag{NSC30}
\]
The extensions are inside $V$ with the cutoff's closed
$U$-support.  The equality near the cancelled part
constructs the actual gluing defining $g_Z$.
The frozen support lemma puts $\operatorname{supp}D\eta$
compactly inside $V$ for every actual test, and full
actual Stokes kills both correction terms.
This is the frozen RAG/RSP interface with a newly
constructed overlap witness, not a general existence
assertion for arbitrary cycles or pivot covers.

\paragraph{Full relations, topology and limits.}
The three-term analytic relation in (NSC9) is the
nonproduct relation itself.  For example $t=1$ and
$(f_i,f_j)=(T_1,T_2)$ retain the entire unit data
$(T_1,T_2,T_1+T_2)$ on its unit locus; the zero-entry
neighborhoods above also handle vanishing proper entries.
For a moving-puncture family one may take
$(f_i,f_j,t)=(u-b,1,b)$, so the transformed entry is $u$,
and append arbitrarily many passive base directions.
All parameter valuations, equalities and mixed
coefficients are retained; no collision extension
or product-probe substitute is inferred.

The cap and triple formulas use only finite actual
scalar operations, full differentiation and the exact
CS/MT operations.  They commute with restrictions and
all analytic refinements.  On each whole fixed compact
presentation, the ordinary zero-gluing cover and the
finite cap refinement give one finite actual output
presentation, with the complete mixed Leibniz seminorms.
No completed form space, current representative,
quotient lifting, global continuous path choice or
arbitrary-direction topology comparison is used.
All norm identities hold at every analytic point over
the original field, including nonrational points;
there is no field-degree renormalization.

ASP's distinct-entry strict obstruction is not contradicted:
the present six-edge mixed relation has zero complete
residual by (NSC8).  No low-degree closed block prohibited
by LTF was used.  The native all-entry scalar versus CNA
comparison remains open, but is no longer needed for
the source-entry result.  General mixed shear-word
relations, arbitrary pivot-cover coherence, full
lower-middle OPEN-KERNEL, general strong G, higher
actual descent and complete BB height remain unproved.
